\documentclass[../main.tex]{subfiles}
\begin{document}

\chapter{Pipelining}
\lessoninfo{
  video={Lesson 3, videos 1--18},
  textbook={H\&P \cite{hennessy2017} Appendix C.1--C.3},
  keywords={pipeline, stage, pipeline CPI, stall, flush, control dependence, data dependence, RAW, WAR, WAW, hazard, forwarding, pipeline depth},
}

Pipelining is the first and most fundamental technique for raising
throughput: instead of finishing one instruction before starting the next,
the processor overlaps the execution of several instructions, each in a
different stage of completion.  This lesson reviews how a pipelined
processor works, quantifies its ideal performance, and then studies the
\emph{dependences} between instructions that prevent the ideal from being
reached.

\section{The idea of pipelining}

Executing an instruction involves a sequence of steps that use different
pieces of hardware.  If the steps are arranged as \term{stage}[ステージ]s
separated by latches, then as soon as instruction $I_1$ leaves stage~1 for
stage~2, instruction $I_2$ can enter stage~1.  With $k$ stages, up to $k$
instructions are in flight at once.\source{Lesson 3, videos 2--3; H\&P
§C.1.}

The consequences for the two metrics of Lesson 2 are different:
\begin{itemize}
  \item \textbf{Throughput} improves by up to a factor of $k$: in steady
    state one instruction completes every stage time instead of every $k$
    stage times.
  \item \textbf{Latency} of an individual instruction does \emph{not}
    improve.  It is the same total work, and the latches between stages add
    a small overhead, so per-instruction latency slightly increases.
\end{itemize}

\subsection{Pipelining in a processor}

The classic five-stage pipeline (\cref{fig:stages}) is used throughout the
course as the reference design.\margin{Real processors have 10--20 stages;
the five-stage model captures every issue that matters here.}

\begin{figure}[htbp]
  \centering
  \input{figures/pipeline-stages}
  \caption{The five-stage instruction pipeline.}
  \label{fig:stages}
\end{figure}

\begin{description}
  \item[IF] \emph{instruction fetch}: read the instruction at the program
    counter (PC) from the instruction memory and increment the PC.
  \item[ID] \emph{instruction decode} and register read: determine the
    operation and read the source operands from the register file.
  \item[EX] \emph{execute}: perform the ALU operation, or compute the
    memory address for a load/store, or evaluate a branch condition.
  \item[MEM] \emph{memory}: access data memory for loads and
    stores.\margin{Each stage is assumed to take exactly one cycle, the
    memory access included---an assumption Lessons 12--15 undo.}
  \item[WB] \emph{write back}: write the result into the register file.
\end{description}

\Cref{fig:timing} shows five independent instructions flowing through the
pipeline.  After the first four cycles fill the pipeline, one instruction
completes per cycle.\caution{Separate instruction and data memory ports are
assumed here; a single port would add \emph{structural} hazards (H\&P
Appendix C).}

\begin{figure}[htbp]
  \centering
  % Five instructions flowing through a five-stage pipeline.
\begin{tikzpicture}[x=0.95cm,y=0.55cm, font=\footnotesize\sffamily,
  st/.style={draw, minimum width=0.9cm, minimum height=0.5cm, inner sep=0pt, fill=ln-box}]
  \foreach \c in {1,...,9} { \node at (\c,0.7) {\c}; }
  \node[anchor=east] at (0.5,0.7) {\iflangja{サイクル}{cycle}};
  \foreach \i/\name in {0/I$_1$,1/I$_2$,2/I$_3$,3/I$_4$,4/I$_5$} {
    \node[anchor=east] at (0.5,-\i) {\name};
    \foreach \s/\lab in {0/IF,1/ID,2/EX,3/MEM,4/WB} {
      \node[st] at (\i+\s+1,-\i) {\lab};
    }
  }
\end{tikzpicture}

  \caption{Instruction timing in a five-stage pipeline: from cycle 5 on, one
    instruction completes every cycle.}
  \label{fig:timing}
\end{figure}

\section{Pipeline CPI}

In the ideal case of \cref{fig:timing} the pipeline completes one
instruction per cycle, so the CPI of the iron law
(\cref{eq:iron-law}) is~1.\source{Lesson 3, video 6; H\&P §C.1.}  Two
effects raise it above~1.
\begin{itemize}
  \item \emph{Fill and drain.}  The first instruction takes $k$ cycles
    before anything completes.  For a program of $n$ instructions the total
    is $k + (n-1)$ cycles, i.e.\ CPI $= (k+n-1)/n \to 1$ as $n$ grows.  For
    real programs ($n$ in the billions) this is negligible.
  \item \emph{Stalls.}  Whenever the pipeline cannot advance an instruction
    to the next stage, a cycle passes without an instruction completing.
    These cycles are what actually determine the CPI of a pipelined
    processor:
    \begin{equation}
      \text{CPI} \;=\; 1 + \frac{\text{stall cycles}}{\text{instruction}} .
      \label{eq:pipeline-cpi}
    \end{equation}
\end{itemize}

\section{Stalls and flushes}

Two mechanisms lose cycles in a pipeline.\source{Lesson 3, videos 7--8;
H\&P §C.2.}

\begin{definition}[Stall]
  A \term{stall}[ストール] holds an instruction (and everything behind it)
  in its current stage for one or more cycles while the instructions ahead
  of it continue.  A bubble---an empty slot---travels down the pipeline in
  its place.
\end{definition}

\begin{definition}[Flush]
  A \term{flush}[フラッシュ] discards instructions that have already been
  fetched into the pipeline because they should not have been executed at
  all, typically because they follow a branch whose outcome was guessed
  wrongly.
\end{definition}

A stall of $c$ cycles costs $c$ cycles; a flush of $m$ instructions costs
$m$ cycles, since each flushed instruction occupied a slot in which a useful
instruction could have completed.

\section{Control dependences}

A \term{control dependence}[制御依存] exists when whether an instruction
executes at all depends on the outcome of an earlier branch.\source{Lesson
3, videos 9--10; H\&P §C.2.}  In the five-stage pipeline the branch
condition and target are known only at the end of EX.\sidenote{Early RISC
pipelines exposed the wasted slots instead of flushing them: in MIPS and
SPARC the instruction after a branch (the \emph{delay slot}) always
executes, and the compiler fills it.  With many instructions in flight one
slot buys little, and MIPS Release~6 (2014) dropped it.}  By then the two
following instructions have already been fetched (\cref{fig:flush}); if the
branch is taken they must be flushed and the fetch restarted at the target.

\begin{figure}[htbp]
  \centering
  % A branch resolved in EX: the two instructions fetched after it are flushed.
\begin{tikzpicture}[x=0.95cm,y=0.55cm, font=\footnotesize\sffamily,
  st/.style={draw, minimum width=0.9cm, minimum height=0.5cm, inner sep=0pt, fill=ln-box},
  bad/.style={st, fill=ln-caution!15, text=ln-caution}]
  \foreach \c in {1,...,8} { \node at (\c,0.7) {\c}; }
  \node[anchor=east] at (0.5,0.7) {\iflangja{サイクル}{cycle}};
  \node[anchor=east] at (0.5,0) {BEQ};
  \foreach \s/\lab in {0/IF,1/ID,2/EX,3/MEM,4/WB} { \node[st] at (\s+1,0) {\lab}; }
  \node[anchor=east] at (0.5,-1) {I$_{\text{next}}$};
  \node[bad] at (2,-1) {IF}; \node[bad] at (3,-1) {ID};
  \node[anchor=east] at (0.5,-2) {I$_{\text{next}+1}$};
  \node[bad] at (3,-2) {IF};
  \node[anchor=east] at (0.5,-3) {\iflangja{分岐先}{target}};
  \foreach \s/\lab in {0/IF,1/ID,2/EX,3/MEM,4/WB} { \node[st] at (\s+4,-3) {\lab}; }
  %% label goes below-LEFT of the arrow tip: centred under it, the (wider)
  %% Japanese word runs into the first box of the branch-target row at x=4.
  \draw[->,ln-caution] (3,0.35) -- (3,-2.4)
    node[below left, inner sep=1pt, font=\scriptsize] {\iflangja{フラッシュ}{flush}};
\end{tikzpicture}

  \caption{A taken branch resolved in EX wastes two fetch slots.}
  \label{fig:flush}
\end{figure}

The cost of this on the CPI depends on how often it
happens.\margin{The \SI{20}{\percent}/\SI{60}{\percent} below are round
lecture numbers; integer code really does branch about every fifth
instruction (H\&P ch.~3).}  Typical
figures used in the lecture are that \SI{20}{\percent} of instructions are
branches and \SI{60}{\percent} of branches are taken; with the simple
policy of always fetching the fall-through instruction (\enquote{predict
not taken}), only taken branches cause a flush:
\begin{equation}
  \text{CPI} = 1 + 0.2 \times 0.6 \times 2 = 1.24 .
  \label{eq:branch-cpi}
\end{equation}

\begin{example}
  In a deeper pipeline where the branch is resolved in stage 10, each taken
  branch flushes nine instructions.  With the same branch statistics the
  CPI becomes $1 + 0.2\times 0.6\times 9 = 2.08$---the processor completes
  less than one instruction every two cycles.  This is why branch prediction
  (Lesson 4) matters more the deeper the pipeline.
\end{example}

\tip{Cycles lost to control dependences $=$ (fraction of instructions that
are branches) $\times$ (fraction mispredicted) $\times$ (instructions
flushed per misprediction).  The last factor is (stage in which the branch
resolves) $-$ 1.}

\section{Data dependences}

A \term{data dependence}[データ依存] exists between two instructions when
they access the same register or memory location and at least one of them
writes it.\source{Lesson 3, videos 11--12; H\&P §C.2.}  There are three
kinds, named by the order of the accesses in program order:

\begin{description}
  \item[RAW (read after write)]%
    \index{read-after-write dependence}%
    the later instruction \emph{reads} a value
    that the earlier one \emph{writes}.  Also called a \emph{true} or
    \emph{flow} dependence: the data really has to flow from one to the
    other.
  \item[WAR (write after read)]%
    \index{write-after-read dependence}%
    the later instruction \emph{writes} a
    location that the earlier one \emph{reads}.  An \emph{anti} dependence.
  \item[WAW (write after write)]%
    \index{write-after-write dependence}%
    both instructions \emph{write} the same
    location; the later write must be the one that survives.  An
    \emph{output} dependence.
\end{description}

WAR and WAW are called \term{false dependence}[偽の依存]s or \emph{name}
dependences: no data moves between the instructions; they merely reuse the
same register name.  They can be removed by renaming (Lesson 6).  A RAW
dependence cannot.\margin{RAR---two reads of the same register---is not a
dependence: order does not matter.}

\begin{example}\leavevmode % head on its own line, listing below it
  \begin{lstlisting}[language={[x86masm]Assembler}, numbers=none]
  I1:  ADD R1, R2, R3     ; R1 <- R2 + R3
  I2:  SUB R4, R1, R5     ; RAW on R1 with I1
  I3:  MUL R1, R6, R7     ; WAW on R1 with I1, WAR with I2
\end{lstlisting}
\end{example}

\section{Dependences and hazards}

A dependence is a property of the \emph{program}; whether it causes a
problem depends on the \emph{pipeline}.\source{Lesson 3, videos 13--14;
H\&P §C.2.}

\begin{definition}[Hazard]
  A \term{hazard}[ハザード] is a situation in which a particular pipeline
  would violate a dependence if it continued unmodified---the dependent
  instruction would read a stale value or execute when it should not.
\end{definition}

In the five-stage pipeline:
\begin{itemize}
  \item A \textbf{RAW} dependence is a hazard only if the instructions are
    close enough.  Registers are written in WB and read in ID, so an
    instruction depends hazardously on the previous one and the one before
    it, but an instruction three or more slots behind reads the register
    after it has been written.\margin{With a register file that writes in
    the first half of the cycle and reads in the second half, the distance
    shrinks by one.}
  \item \textbf{WAR} and \textbf{WAW} dependences never cause hazards here,
    because instructions read registers in ID and write in WB \emph{in
    program order}; a later instruction cannot overtake an earlier one.
    They become hazards in out-of-order pipelines (Lesson 6).
  \item \textbf{Control} dependences are always hazards, because the
    pipeline fetches instructions after the branch before the branch has
    resolved.
\end{itemize}

\begin{keypoint}
  Dependences are decided by the program; hazards by the pipeline.  The
  same program has more hazards in a deeper or out-of-order pipeline.
\end{keypoint}

\section{Handling hazards}

Once a hazard is detected, the pipeline has three options.\source{Lesson 3,
videos 15--16; H\&P §C.2.}

\begin{enumerate}
  \item \textbf{Flush} the dependent instructions---for control hazards,
    where the wrong instructions were fetched.
  \item \textbf{Stall} the dependent instruction until the value it needs
    has been written---always correct, but costs cycles.
  \item \textbf{Fix} the value by \term{forwarding}[フォワーディング]
    (bypassing): route the result directly from the stage that produced it
    to the stage that needs it, without waiting for it to reach the
    register file (\cref{fig:forwarding}).\margin{Lesson 9 adds
    store-to-load forwarding: a pending store hands its value to a later
    load inside the load-store queue.}
\end{enumerate}

\begin{figure}[htbp]
  \centering
  % RAW dependence resolved by forwarding from EX/MEM to EX.
\begin{tikzpicture}[x=0.95cm,y=0.9cm, font=\footnotesize\sffamily,
  st/.style={draw, minimum width=0.9cm, minimum height=0.5cm, inner sep=0pt, fill=ln-box},
  >={Stealth[length=1.6mm]}]
  \foreach \c in {1,...,6} { \node at (\c,0.7) {\c}; }
  \node[anchor=east] at (0.5,0) {\texttt{ADD R1,R2,R3}};
  \foreach \s/\lab in {0/IF,1/ID,2/EX,3/MEM,4/WB} { \node[st] (a\s) at (\s+1,0) {\lab}; }
  \node[anchor=east] at (0.5,-1) {\texttt{SUB R4,R1,R5}};
  \foreach \s/\lab in {0/IF,1/ID,2/EX,3/MEM,4/WB} { \node[st] (b\s) at (\s+2,-1) {\lab}; }
  \draw[->,thick,ln-tip] (a2.south east) to[out=-60,in=120] (b2.north west);
  %% The label sits in the gap between the two rows, to the right of the arc.
  %% Centred at (4.6,-0.35) the (wider) Japanese word prints over the MEM and
  %% WB boxes of the ADD row; y=0.9 opens a gap tall enough to hold it.
  \node[ln-tip, font=\scriptsize, anchor=west] at (3.75,-0.5)
    {\iflangja{フォワーディング}{forwarding}};
\end{tikzpicture}

  \caption{Forwarding: the ALU result of \texttt{ADD} is fed straight into
    the EX stage of \texttt{SUB}, avoiding a stall.}
  \label{fig:forwarding}
\end{figure}

Forwarding removes most RAW stalls but not all of them: a value cannot be
forwarded \emph{backwards in time}.\tip{Stall cycles $=$ (cycle the value
is ready) $-$ (cycle it is needed); $\le 0$ means forwarding alone
suffices.}  The result of a load is only available
at the end of MEM, whereas the next instruction needs it at the start of
its own EX, which is the same cycle.  One stall cycle is therefore
unavoidable for a \emph{load-use} pair even with full forwarding.

\begin{example}\leavevmode % head on its own line, listing below it
  \begin{lstlisting}[language={[x86masm]Assembler}, numbers=none]
  LW  R1, 0(R2)      ; R1 available at end of MEM (cycle 4)
  ADD R3, R1, R4     ; needs R1 in EX (cycle 4): 1 stall
\end{lstlisting}
  With forwarding but no stall the \texttt{ADD} would use a stale R1;
  without forwarding it would stall two cycles.\margin{On a Skylake core an
  L1 hit takes 4 cycles, so a load-use pair there costs about 3 stall
  cycles, not 1 (Intel optimization manual).}
\end{example}

\section{How many stages?}

If pipelining a $k$-stage design gives up to $k\times$ the throughput, why
not use hundreds of stages?\source{Lesson 3, video 17; H\&P §C.1, §1.5.}
Three effects pull in different directions as the depth increases:
\begin{itemize}
  \item \textbf{Clock period decreases} (less logic per stage), which
    lowers $T_{\mathrm{clk}}$ in the iron law---the intended benefit.  But
    the latch overhead per stage does not shrink, so the gain flattens
    out.\caution{A deeper pipeline is not automatically faster: only the
    product CPI $\times T_{\mathrm{clk}}$ decides (Lesson 2).}
  \item \textbf{CPI increases}: more stages means more instructions in
    flight between a branch and its resolution, hence more instructions
    flushed per misprediction, and more distance over which RAW hazards
    need stalls.\margin{A mispredicted branch costs about 16--20 cycles on
    Skylake-class cores (Agner Fog, microarchitecture manual).}
  \item \textbf{Power increases}: more latches switch every cycle and the
    higher frequency raises dynamic power (\cref{eq:active-power}).
\end{itemize}

Considering performance alone (the product CPI $\times T_{\mathrm{clk}}$),
the optimum is deep---on the order of 30--40 stages.  Once power is
included, the optimum moves to roughly 10--15 stages, which is where
current designs sit.\sidenote{The deepest pipeline ever shipped in an x86
was the Pentium~4 \enquote{Prescott} of 2004: 31 stages, up to
\SI{3.8}{\giga\hertz}, with power stopping it there.  Its successor, the
Core microarchitecture of 2006, went back to 14 stages, and depths have
stayed near 15--20 since.}

\begin{keypoint}[Lesson summary]
  Pipelining raises throughput, not latency; ideal CPI is 1 and every stall
  or flush adds to it (\cref{eq:pipeline-cpi}).  Control and RAW dependences
  turn into hazards that cost cycles; forwarding removes most RAW stalls,
  and branch prediction (next lesson) attacks the control hazards.  Pipeline
  depth trades clock period against CPI and power.
\end{keypoint}

\end{document}
